% RESEARCH_DRAFT_SHA256: 6fed5068794baa8e5cdf4d08f1227a2012d7882b46f45a952cf41f10eb749d1f
% Layout and source-reference wrappers expanded for the web renderer.
\documentclass[11pt]{article}
\usepackage[margin=25mm]{geometry}
\usepackage{fontspec,unicode-math}
\setmainfont{Latin Modern Roman}
\setmathfont{Latin Modern Math}
\usepackage{amsmath,amssymb,mathtools,xurl,hyperref}
\hypersetup{colorlinks=true,urlcolor=blue,linkcolor=blue}
\setcounter{secnumdepth}{0}
\setlength{\parindent}{0pt}
\setlength{\parskip}{5pt}
\setlength{\emergencystretch}{3em}
\newcommand{\sref}[2]{\hyperlink{source:#1:#2}{[S#2]}}
\newcommand{\cataloguescope}[1]{\paragraph{Source record.}#1}
\newcommand{\floor}[1]{\left\lfloor#1\right\rfloor}
\begin{document}
\section*{Research-attempt provenance}
This standalone research attempt retains the complete original definition body below,
including its historical source and preparation notes. Those notes describe the
earlier catalogue preparation; the new attempt and its verification are identified
separately after that body. No original ID, catalogue status, or source statement is changed.

\section{Stable Algorithms and Overlap Gaps for Friendly Bisections}\label{statement}

\subsection{Definitions and mathematical statement}

\paragraph{Shared notation.}Standard number systems are $\mathbb N,\mathbb Z,\mathbb Q,\mathbb R,\mathbb C$. Specialized notation and index ranges follow the cited source and the definitions below.


Polynomial time means a number of elementary steps bounded by a polynomial in the encoded input length. A decision problem is in $\mathrm{NP}$ if yes-instances have polynomial-length certificates verifiable in polynomial time; NP-hardness is under polynomial-time reductions.

A bisection partitions the vertices into two parts of equal size (or differing by one for odd order). Its cut consists of the edges joining the parts; minimum bisection minimizes their total weight.
\paragraph{Record scope and qualifications.}
The catalogue title is retained as metadata. Named certificates, contracts, bounds or frameworks in that title are not asserted as results of the original source extract.
Source-record limitation: The source-specific h-friendly threshold and stable-algorithm model are absent from the cleaned statement; inventing their normalizations would change the question.
\paragraph{Statement recorded in the supplied source.}
For $h>0$, is there a stable polynomial-time algorithm that finds an $h$-friendly bisection of $G(n,1/2)$? Does the overlap-gap property hold for every $h>0$, perhaps by transfer from random regular graphs to Erdős–Rényi graphs? \href{https://doi.org/10.4171/owr/2023/1}{[S2]}
\paragraph{Source status.}
The supplied catalogue labels this source record open; that metadata was checked against the supplied snapshot on 2026-10-11. The label is preserved for provenance and is not an independent verification of present open status or mathematical correctness. The catalogue literature assessment is dated 2026-08-22.
\subsection{Short English statement}
The source asks about stable Algorithms and Overlap Gaps for Friendly Bisections.
\subsection{Sources}
\begin{itemize}
\item[S1] Ulam AI and the UnsolvedMath Contributors, UnsolvedMath: A Curated Collection of Open Mathematics Problems, catalogue entry. Catalogue curation and metadata: CC BY 4.0. The source extract is retained; definitions, English summaries and qualifications are editorial additions.
\url{https://huggingface.co/datasets/ulamai/UnsolvedMath}.
\item[S2] Combinatorics (2023), Oberwolfach Reports, DOI 10.4171/owr/2023/1, pp. 5--89
\url{https://doi.org/10.4171/owr/2023/1}.
\item[S3] FKN, Friendly Bisections in Random Graphs
\url{https://arxiv.org/abs/2105.13337}.
\item[S4] Overlap Gaps in Random Regular Graphs
\url{https://arxiv.org/abs/2202.10379}.
\item Source-specific Oberwolfach Reports licences, including CC BY-SA 4.0 where applicable, are retained; see THIRD\_PARTY\_NOTICES.md.
\end{itemize}
\paragraph{Preparation and review.}
The original source and UnsolvedMath contributors retain attribution for the source extract. Codex (GPT-6; exact serving revision not exposed) assisted with definitions, English summaries, scope notes and formatting. No new human mathematical authorship or independently verified solution is asserted. Catalogue identity checked against the complete supplied registry and current repository. Distinct upstream records, including repeated questions and broad or damaged extracts, are retained for complete coverage of the source CS category. A repeated record is not a new independent conjecture. Complete mathematical-equivalence and current-literature review has not been performed. Preparation and final review follow contribution prompts C and D, with source gaps disclosed. Structural and rendering checks do not establish a new theorem or open status.
\label{end:statement}

\section{Literature and context check}

I read Gamarnik's primary contribution on pages 61--63 of the 2023 report.
It defines the threshold as $(d_{\rm same}-d_{\rm opposite})/\sqrt n\geq h$
and distinguishes fully friendly from asymptotically friendly partitions.
The report discusses an overlap gap for a stated interval of positive $h$ and
leaves the broader transfer question open. The calculations below use this
normalization, but do not define an absent universal model of ``stable algorithm''
or claim a barrier for every polynomial-time method.

\section{Research attempt: Swap optimality, spectral necessity, and failure of an energy-only transfer}

\paragraph{What local minimum bisection actually guarantees.}
Let $n$ be even, let $A$ be the adjacency matrix of a finite simple graph,
and let $\sigma_i\in\{-1,1\}$ encode an exactly balanced partition.
Put $D_i=\sum_jA_{ij}\sigma_i\sigma_j$, the same-side degree minus the
opposite-side degree. Swap opposite-side vertices $u,v$.
Flipping one vertex changes the cut size by its $D_i$. Adding the two
single-flip changes counts their shared crossing edge as a decrease twice,
although after the simultaneous swap it still crosses. Thus the true change is
\[
 \Delta\operatorname{cut}=D_u+D_v+2A_{uv}.
\]
A partition locally minimizing cut under swaps therefore gives only
$D_u+D_v+2A_{uv}\geq0$ for opposite-side pairs. It does not give
$D_i\geq0$ separately, still less $D_i\geq h\sqrt n$.
For an explicit example, consider the four-vertex star, with its center and
one leaf on one side. The center has $D=-1$, its same-side leaf has $D=1$,
and the other two leaves have $D=-1$. Every bisection of this graph cuts two
edges, so this is a global minimum bisection, yet three vertices are unfriendly.
Equivalently the swap formula is zero for every opposite-side pair.
This does not refute a random-graph theorem; it refutes a deterministic inference
used in a proposed greedy route.

\paragraph{A necessary spectral inequality, with the centering term retained.}
Write $J$ for the all-ones matrix and
$W=A-\frac12(J-I)$. Balancedness gives $\sigma^TJ\sigma=0$ and
$\sigma^T\sigma=n$, hence
\[
 \sum_iD_i=\sigma^TA\sigma=\sigma^TW\sigma-\frac n2
 \leq n\|W\|_{\rm op}-\frac n2.
\]
A fully $h$-friendly bisection must therefore satisfy
\[
 h\leq\frac{\|W\|_{\rm op}}{\sqrt n}-\frac1{2\sqrt n}.
\]
This is a finite-instance necessary condition only. A bound on the sum of local
fields is not a lower bound on their minimum, and it says nothing about the
pairwise geometry of the solution set needed for an overlap gap.

\paragraph{Why an energy transfer alone does not transfer friendliness.}
Changing $q$ undirected edges changes $\sigma^TA\sigma$ by at most $2q$
for a fixed $\sigma$. With $q\leq n$, its normalization by $n\sqrt n$
changes by at most $2/\sqrt n$. But changes incident to one vertex can alter
that vertex's normalized local field by order $\sqrt n$.
For an explicit family, take two disjoint cliques of size $m\geq2$, partitioned
by their components, so $n=2m$. Choose one vertex $u$ in the first clique,
remove its $m-1$ edges within that clique, and add its $m$ edges to the
other clique. Exactly $2m-1<n$ edges change. Its field changes from
$m-1$ to $-m$, a difference $2m-1$, while the change in the normalized total
energy is at most $2/\sqrt n$. Every other vertex has a nonnegative field
for $m\geq2$, but $u$ violates every positive friendliness threshold.
Thus a proof comparing only normalized total energies can lose a vertexwise
constraint even when its energy error vanishes. For an asymptotic criterion
allowing $o(n)$ bad vertices this one-vertex example is admissible instead;
it must not be used as a counterexample to that weaker criterion.

\section{Verification and adversarial audit}

The swap correction has a plus sign and is checked on all swaps of the
four-vertex star. The centered spectral identity includes the diagonal term
$-n/2$. The transfer example counts fewer than $n$ changed edges and states
that a single exceptional vertex does not obstruct asymptotic friendliness.
No deterministic construction is identified with a typical $G(n,1/2)$ sample.

\paragraph{Audit corrections.} Recovered the source threshold normalization and separated full from
asymptotic friendliness. Kept the energy example from being promoted to an
asymptotic or random-graph impossibility theorem.
\paragraph{External computation.} None; symbolic arguments and boundary checks only.
\section{Scope and remaining gaps}

The route eliminates two unjustified sufficient conditions for full friendliness
and gives a finite-instance necessary bound. It proves neither a stable random-
graph algorithm nor an overlap gap for every positive $h$, and gives no transfer
of an algorithmic barrier between random regular and Erd\H{o}s--R\'enyi models.
The original stability and asymptotic quantifiers remain essential gaps.

\paragraph{Integration review.} Codex (GPT-6; exact serving revision not exposed) reviewed the written arguments and expanded layout and source-reference wrappers. Repository metadata and the production web renderer were checked using Python 3.13 and JavaScript with MathJax 3.2.2. This model-assisted review is not independent human proof verification; no PDF compilation or formal proof checking is asserted.

\section{Completion Estimate}
\noindent\textbf{COMPLETION: 0\%}
The calculations below the preserved source statement investigate a concrete route,
but establish no substantiated new progress toward the entire original target.
Known special cases and failed sufficient conditions are not credited as a solution.
\section{References}
The original source references and catalogue attribution are preserved in the
Sources section above. The following URLs were actually consulted for this attempt;
retention of an additional original reference does not claim fresh consultation of it.
\begin{itemize}
\item \url{https://ems.press/content/serial-article-files/46994} (consulted 2026-10-11). 
\end{itemize}
\section{Model and contributor disclosure}
Codex GPT-6 completed written-proof review of this attempt at its declared scope; this is model-assisted review, not independent human proof checking.
OpenAI gpt-6-astra, requested reasoning configuration ultra, generated this research
attempt on 2026-10-11. The exact serving revision is not exposed. Codex GPT-6
performed orchestration and repository-format review; its exact serving revision
is also not exposed. The mathematical route, boundary checks and final adversarial
self-audit followed contribution prompts A then B in
\texttt{docs/contribution\_prompts.md}, after reading \texttt{CONTRIBUTING.md}.
Python 3.14.2 assembled the standalone documents; any mathematical computation is
specified in Verification. No theorem prover was used. No human mathematical
authorship, substantial human proof correction, or independent human proof checking
is asserted. Model-assisted review, source retrieval, and structural checks do not
establish correctness or novelty. Full private reasoning transcripts are unavailable
as reproducibility artifacts; the complete arguments, inputs and outcomes actually
used are stated here or in the companion batch audit.
\section{Final status}
Unresolved as a new research contribution. This is a concrete route investigation,
not a claim of new complete resolution or a current-status certification. A known
answer reproduced for context does not become a new contribution. The catalogue's
source-status label is preserved; this attempt status does not assert that a
source-labelled solved question is open.
\end{document}
